% Môn: Vật lý
% Lớp: 12
% Chương: Chương II. Khí lí tưởng
% Bài: L12C2 Bài 11. Phương trình trạng thái của khí lí tưởng · Các thông số trạng thái của một lượng khí
% ===== Câu 3 =====
% ID: L12C2B11-87-TLN
% Mức: TH
\dangbt{Các thông số trạng thái của một lượng khí}
\begin{ex}
% ID: L12C2B11-87-TLN
% Mức: TH
Trong dạng “Chuyển động Brown và bằng chứng về chuyển động phân tử”, Một bình chứa $6\times10^{23}$ phân tử trong thể tích $0,02\,m^3$. Tính mật độ phân tử theo đơn vị $10^{25}\,m^{-3}$.
\shortans{3}
\loigiai{
$n_0=N/V=6\times10^{23}/0,02=3\times10^{25}\,m^{-3}$.
}
\end{ex}

% ===== Câu 7 =====
% ID: L12C2B11-88-TLN
% Mức: TH
\begin{ex}
% ID: L12C2B11-88-TLN
% Mức: TH
Trong dạng “Chuyển động Brown và bằng chứng về chuyển động phân tử”, Một lượng khí có thể tích 4 lít được nén còn 2 lít. Mật độ phân tử tăng bao nhiêu lần?
\shortans{2}
\loigiai{
Số phân tử không đổi nên mật độ tỉ lệ nghịch thể tích: $n_2/n_1=V_1/V_2=2$.
}
\end{ex}

% ===== Câu 15 =====
% ID: L12C2B11-89-TLN
% Mức: VD
\begin{ex}
% ID: L12C2B11-89-TLN
% Mức: VD
Trong dạng “Chuyển động Brown và bằng chứng về chuyển động phân tử”, Khí có $N=1,2\times10^{24}$ phân tử và mật độ $3\times10^{25}\,m^{-3}$. Tính thể tích theo lít.
\shortans{40}
\loigiai{
$V=N/n_0=0,04\,m^3=40$ lít.
}
\end{ex}

% ===== Câu 19 =====
% ID: L12C2B11-90-TLN
% Mức: VD
\begin{ex}
% ID: L12C2B11-90-TLN
% Mức: VD
Trong dạng “Chuyển động Brown và bằng chứng về chuyển động phân tử”, Có 2 mol khí. Lấy $N_A=6,02\times10^{23}$. Số phân tử theo đơn vị $10^{23}$ là bao nhiêu?
\shortans{12,04}
\loigiai{
$N=2N_A=12,04\times10^{23}$.
}
\end{ex}

% ===== Câu 23 =====
% ID: L12C2B11-91-TLN
% Mức: VDC
\begin{ex}
% ID: L12C2B11-91-TLN
% Mức: VDC
Trong dạng “Chuyển động Brown và bằng chứng về chuyển động phân tử”, Khí chiếm 5 lít, giãn nở thành 15 lít. Mật độ phân tử giảm bao nhiêu lần?
\shortans{3}
\loigiai{
$n_1/n_2=V_2/V_1=15/5=3$.
}
\end{ex}

% ===== Câu 27 =====
% ID: L12C2B11-92-TLN
% Mức: TH
\begin{ex}
% ID: L12C2B11-92-TLN
% Mức: TH
Trong dạng “Các luận điểm cơ bản của mô hình động học phân tử chất khí”, Một bình chứa $6\times10^{23}$ phân tử trong thể tích $0,02\,m^3$. Tính mật độ phân tử theo đơn vị $10^{25}\,m^{-3}$.
\shortans{3}
\loigiai{
$n_0=N/V=6\times10^{23}/0,02=3\times10^{25}\,m^{-3}$.
}
\end{ex}

% ===== Câu 28 =====
% ID: L12C2B11-93-TN
% Mức: VD
\begin{ex}
% ID: L12C2B11-93-TN
% Mức: VD
Trong dạng “Chuyển động Brown và bằng chứng về chuyển động phân tử”, Ba thông số trạng thái cơ bản của một lượng khí là
\choice
{\True áp suất, thể tích, nhiệt độ}
{khối lượng, trọng lượng, nhiệt lượng}
{vận tốc, gia tốc, thời gian}
{lực, công, công suất}
\loigiai{
Đáp án A. Trạng thái lượng khí xác định bởi p, $V$, T.
}
\end{ex}

% ===== Câu 31 =====
% ID: L12C2B11-94-TLN
% Mức: TH
\begin{ex}
% ID: L12C2B11-94-TLN
% Mức: TH
Trong dạng “Các luận điểm cơ bản của mô hình động học phân tử chất khí”, Một lượng khí có thể tích 4 lít được nén còn 2 lít. Mật độ phân tử tăng bao nhiêu lần?
\shortans{2}
\loigiai{
Số phân tử không đổi nên mật độ tỉ lệ nghịch thể tích: $n_2/n_1=V_1/V_2=2$.
}
\end{ex}

% ===== Câu 36 =====
% ID: L12C2B11-95-TN
% Mức: VDC
\begin{ex}
% ID: L12C2B11-95-TN
% Mức: VDC
Trong dạng “Chuyển động Brown và bằng chứng về chuyển động phân tử”, Khi nén khí ở nhiệt độ không đổi, mật độ phân tử
\choice
{\True tăng}
{giảm}
{không đổi}
{bằng không}
\loigiai{
Đáp án A. Cùng số phân tử trong thể tích nhỏ hơn nên mật độ phân tử tăng.
}
\end{ex}

% ===== Câu 39 =====
% ID: L12C2B11-96-TLN
% Mức: VD
\begin{ex}
% ID: L12C2B11-96-TLN
% Mức: VD
Trong dạng “Các luận điểm cơ bản của mô hình động học phân tử chất khí”, Khí có $N=1,2\times10^{24}$ phân tử và mật độ $3\times10^{25}\,m^{-3}$. Tính thể tích theo lít.
\shortans{40}
\loigiai{
$V=N/n_0=0,04\,m^3=40$ lít.
}
\end{ex}

% ===== Câu 41 =====
% ID: L12C2B11-97-DS
% Mức: TH
\begin{ex}
% ID: L12C2B11-97-DS
% Mức: TH
Xét các phát biểu sau trong dạng “Các thông số trạng thái của một lượng khí”.
\choiceTF
{\True Theo mô hình động học phân tử, các phân tử khí — phát biểu đúng là: chuyển động hỗn loạn không ngừng.}
{Theo mô hình động học phân tử, các phân tử khí — phát biểu đứng yên tại vị trí cân bằng là đúng.}
{\True Kích thước phân tử so với khoảng cách trung bình giữa chúng trong khí lí tưởng được coi là — phát biểu đúng là: rất nhỏ.}
{Kích thước phân tử so với khoảng cách trung bình giữa chúng trong khí lí tưởng được coi là — phát biểu bằng nhau là đúng.}
\loigiai{
\begin{itemchoice}
\item (Đúng) Theo mô hình động học phân tử, các phân tử khí — phát biểu đúng là: chuyển động hỗn loạn không ngừng. (Vì): Phân tử khí luôn chuyển động nhiệt hỗn loạn không ngừng. b) Sai: nội dung không phù hợp. c) Đúng: Trong mô hình khí lí tưởng, kích thước phân tử được bỏ qua so với khoảng cách giữa chúng. d) Sai: nội dung không phù hợp
\item (Sai) Theo mô hình động học phân tử, các phân tử khí — phát biểu đứng yên tại vị trí cân bằng là đúng.
\item (Đúng) Kích thước phân tử so với khoảng cách trung bình giữa chúng trong khí lí tưởng được coi là — phát biểu đúng là: rất nhỏ.
\item (Sai) Kích thước phân tử so với khoảng cách trung bình giữa chúng trong khí lí tưởng được coi là — phát biểu bằng nhau là đúng.
\end{itemchoice}
}
\end{ex}

% ===== Câu 43 =====
% ID: L12C2B11-98-TLN
% Mức: VD
\begin{ex}
% ID: L12C2B11-98-TLN
% Mức: VD
Trong dạng “Các luận điểm cơ bản của mô hình động học phân tử chất khí”, Có 2 mol khí. Lấy $N_A=6,02\times10^{23}$. Số phân tử theo đơn vị $10^{23}$ là bao nhiêu?
\shortans{12,04}
\loigiai{
$N=2N_A=12,04\times10^{23}$.
}
\end{ex}

% ===== Câu 45 =====
% ID: L12C2B11-99-DS
% Mức: TH
\begin{ex}
% ID: L12C2B11-99-DS
% Mức: TH
Xét các phát biểu sau trong dạng “Các thông số trạng thái của một lượng khí”.
\choiceTF
{\True Lực tương tác giữa các phân tử khí lí tưởng được coi là — phát biểu đúng là: không đáng kể trừ khi va chạm.}
{Lực tương tác giữa các phân tử khí lí tưởng được coi là — phát biểu luôn rất lớn là đúng.}
{\True Chuyển động Brown là chuyển động — phát biểu đúng là: hỗn loạn của hạt nhỏ lơ lửng do va chạm phân tử.}
{Chuyển động Brown là chuyển động — phát biểu rơi tự do của hạt bụi là đúng.}
\loigiai{
\begin{itemchoice}
\item (Đúng) Lực tương tác giữa các phân tử khí lí tưởng được coi là — phát biểu đúng là: không đáng kể trừ khi va chạm. (Vì): Ngoài thời gian va chạm rất ngắn, tương tác phân tử được bỏ qua. b) Sai: nội dung không phù hợp. c) Đúng: Các phân tử môi trường va chạm không cân bằng làm hạt nhỏ chuyển động Brown. d) Sai: nội dung không phù hợp
\item (Sai) Lực tương tác giữa các phân tử khí lí tưởng được coi là — phát biểu luôn rất lớn là đúng.
\item (Đúng) Chuyển động Brown là chuyển động — phát biểu đúng là: hỗn loạn của hạt nhỏ lơ lửng do va chạm phân tử.
\item (Sai) Chuyển động Brown là chuyển động — phát biểu rơi tự do của hạt bụi là đúng.
\end{itemchoice}
}
\end{ex}

% ===== Câu 47 =====
% ID: L12C2B11-100-TLN
% Mức: VDC
\begin{ex}
% ID: L12C2B11-100-TLN
% Mức: VDC
Trong dạng “Các luận điểm cơ bản của mô hình động học phân tử chất khí”, Khí chiếm 5 lít, giãn nở thành 15 lít. Mật độ phân tử giảm bao nhiêu lần?
\shortans{3}
\loigiai{
$n_1/n_2=V_2/V_1=15/5=3$.
}
\end{ex}

% ===== Câu 49 =====
% ID: L12C2B11-101-DS
% Mức: VD
\begin{ex}
% ID: L12C2B11-101-DS
% Mức: VD
Xét các phát biểu sau trong dạng “Các thông số trạng thái của một lượng khí”.
\choiceTF
{\True Áp suất khí lên thành bình sinh ra chủ yếu do — phát biểu đúng là: phân tử va chạm thành bình.}
{Áp suất khí lên thành bình sinh ra chủ yếu do — phát biểu trọng lượng bình là đúng.}
{\True Khi nhiệt độ khí tăng, chuyển động nhiệt của phân tử nhìn chung — phát biểu đúng là: nhanh hơn.}
{Khi nhiệt độ khí tăng, chuyển động nhiệt của phân tử nhìn chung — phát biểu không đổi là đúng.}
\loigiai{
\begin{itemchoice}
\item (Đúng) Áp suất khí lên thành bình sinh ra chủ yếu do — phát biểu đúng là: phân tử va chạm thành bình. (Vì): Xung lượng truyền trong các va chạm phân tử với thành bình tạo áp suất. b) Sai: nội dung không phù hợp. c) Đúng: Nhiệt độ tăng làm động năng tịnh tiến trung bình tăng. d) Sai: nội dung không phù hợp
\item (Sai) Áp suất khí lên thành bình sinh ra chủ yếu do — phát biểu trọng lượng bình là đúng.
\item (Đúng) Khi nhiệt độ khí tăng, chuyển động nhiệt của phân tử nhìn chung — phát biểu đúng là: nhanh hơn.
\item (Sai) Khi nhiệt độ khí tăng, chuyển động nhiệt của phân tử nhìn chung — phát biểu không đổi là đúng.
\end{itemchoice}
}
\end{ex}

% ===== Câu 51 =====
% ID: L12C2B11-102-TLN
% Mức: TH
\begin{ex}
% ID: L12C2B11-102-TLN
% Mức: TH
Trong dạng “Các thông số trạng thái của một lượng khí”, Một bình chứa $6\times10^{23}$ phân tử trong thể tích $0,02\,m^3$. Tính mật độ phân tử theo đơn vị $10^{25}\,m^{-3}$.
\shortans{3}
\loigiai{
$n_0=N/V=6\times10^{23}/0,02=3\times10^{25}\,m^{-3}$.
}
\end{ex}

% ===== Câu 53 =====
% ID: L12C2B11-103-DS
% Mức: VD
\begin{ex}
% ID: L12C2B11-103-DS
% Mức: VD
Xét các phát biểu sau trong dạng “Các thông số trạng thái của một lượng khí”.
\choiceTF
{\True Ba thông số trạng thái cơ bản của một lượng khí là — phát biểu đúng là: áp suất, thể tích, nhiệt độ.}
{Ba thông số trạng thái cơ bản của một lượng khí là — phát biểu khối lượng, trọng lượng, nhiệt lượng là đúng.}
{\True Va chạm giữa các phân tử khí lí tưởng được xem là — phát biểu đúng là: đàn hồi.}
{Va chạm giữa các phân tử khí lí tưởng được xem là — phát biểu không bảo toàn động lượng là đúng.}
\loigiai{
\begin{itemchoice}
\item (Đúng) Ba thông số trạng thái cơ bản của một lượng khí là — phát biểu đúng là: áp suất, thể tích, nhiệt độ. (Vì): Trạng thái lượng khí xác định bởi p, $V$, T. b) Sai: nội dung không phù hợp. c) Đúng: Mô hình khí lí tưởng coi các va chạm là đàn hồi. d) Sai: nội dung không phù hợp
\item (Sai) Ba thông số trạng thái cơ bản của một lượng khí là — phát biểu khối lượng, trọng lượng, nhiệt lượng là đúng.
\item (Đúng) Va chạm giữa các phân tử khí lí tưởng được xem là — phát biểu đúng là: đàn hồi.
\item (Sai) Va chạm giữa các phân tử khí lí tưởng được xem là — phát biểu không bảo toàn động lượng là đúng.
\end{itemchoice}
}
\end{ex}

% ===== Câu 55 =====
% ID: L12C2B11-104-TLN
% Mức: TH
\begin{ex}
% ID: L12C2B11-104-TLN
% Mức: TH
Trong dạng “Các thông số trạng thái của một lượng khí”, Một lượng khí có thể tích 4 lít được nén còn 2 lít. Mật độ phân tử tăng bao nhiêu lần?
\shortans{2}
\loigiai{
Số phân tử không đổi nên mật độ tỉ lệ nghịch thể tích: $n_2/n_1=V_1/V_2=2$.
}
\end{ex}

% ===== Câu 57 =====
% ID: L12C2B11-105-DS
% Mức: VDC
\begin{ex}
% ID: L12C2B11-105-DS
% Mức: VDC
Xét các phát biểu sau trong dạng “Các thông số trạng thái của một lượng khí”.
\choiceTF
{\True Khi nén khí ở nhiệt độ không đổi, mật độ phân tử — phát biểu đúng là: tăng.}
{Khi nén khí ở nhiệt độ không đổi, mật độ phân tử — phát biểu giảm là đúng.}
{\True Hiện tượng khuếch tán chứng tỏ các phân tử — phát biểu đúng là: chuyển động không ngừng.}
{Hiện tượng khuếch tán chứng tỏ các phân tử — phát biểu đứng yên là đúng.}
\loigiai{
\begin{itemchoice}
\item (Đúng) Khi nén khí ở nhiệt độ không đổi, mật độ phân tử — phát biểu đúng là: tăng. (Vì): Cùng số phân tử trong thể tích nhỏ hơn nên mật độ phân tử tăng. b) Sai: nội dung không phù hợp. c) Đúng: Khuếch tán là hệ quả trực tiếp của chuyển động nhiệt hỗn loạn. d) Sai: nội dung không phù hợp
\item (Sai) Khi nén khí ở nhiệt độ không đổi, mật độ phân tử — phát biểu giảm là đúng.
\item (Đúng) Hiện tượng khuếch tán chứng tỏ các phân tử — phát biểu đúng là: chuyển động không ngừng.
\item (Sai) Hiện tượng khuếch tán chứng tỏ các phân tử — phát biểu đứng yên là đúng.
\end{itemchoice}
}
\end{ex}

% ===== Câu 63 =====
% ID: L12C2B11-106-TLN
% Mức: VD
\begin{ex}
% ID: L12C2B11-106-TLN
% Mức: VD
Trong dạng “Các thông số trạng thái của một lượng khí”, Khí có $N=1,2\times10^{24}$ phân tử và mật độ $3\times10^{25}\,m^{-3}$. Tính thể tích theo lít.
\shortans{40}
\loigiai{
$V=N/n_0=0,04\,m^3=40$ lít.
}
\end{ex}

% ===== Câu 67 =====
% ID: L12C2B11-107-TLN
% Mức: VD
\begin{ex}
% ID: L12C2B11-107-TLN
% Mức: VD
Trong dạng “Các thông số trạng thái của một lượng khí”, Có 2 mol khí. Lấy $N_A=6,02\times10^{23}$. Số phân tử theo đơn vị $10^{23}$ là bao nhiêu?
\shortans{12,04}
\loigiai{
$N=2N_A=12,04\times10^{23}$.
}
\end{ex}

% ===== Câu 68 =====
% ID: L12C2B11-108-TN
% Mức: VD
\begin{ex}
% ID: L12C2B11-108-TN
% Mức: VD
Trong dạng “Các luận điểm cơ bản của mô hình động học phân tử chất khí”, Ba thông số trạng thái cơ bản của một lượng khí là
\choice
{\True áp suất, thể tích, nhiệt độ}
{khối lượng, trọng lượng, nhiệt lượng}
{vận tốc, gia tốc, thời gian}
{lực, công, công suất}
\loigiai{
Đáp án A. Trạng thái lượng khí xác định bởi p, $V$, T.
}
\end{ex}

% ===== Câu 71 =====
% ID: L12C2B11-109-TLN
% Mức: VDC
\begin{ex}
% ID: L12C2B11-109-TLN
% Mức: VDC
Trong dạng “Các thông số trạng thái của một lượng khí”, Khí chiếm 5 lít, giãn nở thành 15 lít. Mật độ phân tử giảm bao nhiêu lần?
\shortans{3}
\loigiai{
$n_1/n_2=V_2/V_1=15/5=3$.
}
\end{ex}

% ===== Câu 75 =====
% ID: L12C2B11-110-TLN
% Mức: TH
\begin{ex}
% ID: L12C2B11-110-TLN
% Mức: TH
Trong dạng “Giải thích áp suất khí bằng va chạm phân tử”, Một bình chứa $6\times10^{23}$ phân tử trong thể tích $0,02\,m^3$. Tính mật độ phân tử theo đơn vị $10^{25}\,m^{-3}$.
\shortans{3}
\loigiai{
$n_0=N/V=6\times10^{23}/0,02=3\times10^{25}\,m^{-3}$.
}
\end{ex}

% ===== Câu 76 =====
% ID: L12C2B11-111-TN
% Mức: VDC
\begin{ex}
% ID: L12C2B11-111-TN
% Mức: VDC
Trong dạng “Các luận điểm cơ bản của mô hình động học phân tử chất khí”, Khi nén khí ở nhiệt độ không đổi, mật độ phân tử
\choice
{\True tăng}
{giảm}
{không đổi}
{bằng không}
\loigiai{
Đáp án A. Cùng số phân tử trong thể tích nhỏ hơn nên mật độ phân tử tăng.
}
\end{ex}

% ===== Câu 79 =====
% ID: L12C2B11-112-TLN
% Mức: TH
\begin{ex}
% ID: L12C2B11-112-TLN
% Mức: TH
Trong dạng “Giải thích áp suất khí bằng va chạm phân tử”, Một lượng khí có thể tích 4 lít được nén còn 2 lít. Mật độ phân tử tăng bao nhiêu lần?
\shortans{2}
\loigiai{
Số phân tử không đổi nên mật độ tỉ lệ nghịch thể tích: $n_2/n_1=V_1/V_2=2$.
}
\end{ex}

% ===== Câu 87 =====
% ID: L12C2B11-113-TLN
% Mức: VD
\begin{ex}
% ID: L12C2B11-113-TLN
% Mức: VD
Trong dạng “Giải thích áp suất khí bằng va chạm phân tử”, Khí có $N=1,2\times10^{24}$ phân tử và mật độ $3\times10^{25}\,m^{-3}$. Tính thể tích theo lít.
\shortans{40}
\loigiai{
$V=N/n_0=0,04\,m^3=40$ lít.
}
\end{ex}

% ===== Câu 91 =====
% ID: L12C2B11-114-TLN
% Mức: VD
\begin{ex}
% ID: L12C2B11-114-TLN
% Mức: VD
Trong dạng “Giải thích áp suất khí bằng va chạm phân tử”, Có 2 mol khí. Lấy $N_A=6,02\times10^{23}$. Số phân tử theo đơn vị $10^{23}$ là bao nhiêu?
\shortans{12,04}
\loigiai{
$N=2N_A=12,04\times10^{23}$.
}
\end{ex}

% ===== Câu 95 =====
% ID: L12C2B11-115-TLN
% Mức: VDC
\begin{ex}
% ID: L12C2B11-115-TLN
% Mức: VDC
Trong dạng “Giải thích áp suất khí bằng va chạm phân tử”, Khí chiếm 5 lít, giãn nở thành 15 lít. Mật độ phân tử giảm bao nhiêu lần?
\shortans{3}
\loigiai{
$n_1/n_2=V_2/V_1=15/5=3$.
}
\end{ex}

% ===== Câu 99 =====
% ID: L12C2B11-116-TLN
% Mức: TH
\begin{ex}
% ID: L12C2B11-116-TLN
% Mức: TH
Trong dạng “Vận dụng mô hình động học phân tử vào thực tế”, Một bình chứa $6\times10^{23}$ phân tử trong thể tích $0,02\,m^3$. Tính mật độ phân tử theo đơn vị $10^{25}\,m^{-3}$.
\shortans{3}
\loigiai{
$n_0=N/V=6\times10^{23}/0,02=3\times10^{25}\,m^{-3}$.
}
\end{ex}

% ===== Câu 103 =====
% ID: L12C2B11-117-TLN
% Mức: TH
\begin{ex}
% ID: L12C2B11-117-TLN
% Mức: TH
Trong dạng “Vận dụng mô hình động học phân tử vào thực tế”, Một lượng khí có thể tích 4 lít được nén còn 2 lít. Mật độ phân tử tăng bao nhiêu lần?
\shortans{2}
\loigiai{
Số phân tử không đổi nên mật độ tỉ lệ nghịch thể tích: $n_2/n_1=V_1/V_2=2$.
}
\end{ex}

% ===== Câu 108 =====
% ID: L12C2B11-118-TN
% Mức: VD
\begin{ex}
% ID: L12C2B11-118-TN
% Mức: VD
Trong dạng “Các thông số trạng thái của một lượng khí”, Ba thông số trạng thái cơ bản của một lượng khí là
\choice
{\True áp suất, thể tích, nhiệt độ}
{khối lượng, trọng lượng, nhiệt lượng}
{vận tốc, gia tốc, thời gian}
{lực, công, công suất}
\loigiai{
Đáp án A. Trạng thái lượng khí xác định bởi p, $V$, T.
}
\end{ex}

% ===== Câu 111 =====
% ID: L12C2B11-119-TLN
% Mức: VD
\begin{ex}
% ID: L12C2B11-119-TLN
% Mức: VD
Trong dạng “Vận dụng mô hình động học phân tử vào thực tế”, Khí có $N=1,2\times10^{24}$ phân tử và mật độ $3\times10^{25}\,m^{-3}$. Tính thể tích theo lít.
\shortans{40}
\loigiai{
$V=N/n_0=0,04\,m^3=40$ lít.
}
\end{ex}

% ===== Câu 115 =====
% ID: L12C2B11-120-TLN
% Mức: VD
\begin{ex}
% ID: L12C2B11-120-TLN
% Mức: VD
Trong dạng “Vận dụng mô hình động học phân tử vào thực tế”, Có 2 mol khí. Lấy $N_A=6,02\times10^{23}$. Số phân tử theo đơn vị $10^{23}$ là bao nhiêu?
\shortans{12,04}
\loigiai{
$N=2N_A=12,04\times10^{23}$.
}
\end{ex}

% ===== Câu 116 =====
% ID: L12C2B11-121-TN
% Mức: VDC
\begin{ex}
% ID: L12C2B11-121-TN
% Mức: VDC
Trong dạng “Các thông số trạng thái của một lượng khí”, Khi nén khí ở nhiệt độ không đổi, mật độ phân tử
\choice
{\True tăng}
{giảm}
{không đổi}
{bằng không}
\loigiai{
Đáp án A. Cùng số phân tử trong thể tích nhỏ hơn nên mật độ phân tử tăng.
}
\end{ex}

% ===== Câu 119 =====
% ID: L12C2B11-122-TLN
% Mức: VDC
\begin{ex}
% ID: L12C2B11-122-TLN
% Mức: VDC
Trong dạng “Vận dụng mô hình động học phân tử vào thực tế”, Khí chiếm 5 lít, giãn nở thành 15 lít. Mật độ phân tử giảm bao nhiêu lần?
\shortans{3}
\loigiai{
$n_1/n_2=V_2/V_1=15/5=3$.
}
\end{ex}

% ===== Câu 122 =====
% ID: L12C2B11-123-TLN
% Mức: TH
\begin{ex}
% ID: L12C2B11-123-TLN
% Mức: TH
Trong dạng “Đặc điểm cấu tạo và chuyển động phân tử ở thể khí”, Một bình chứa $6\times10^{23}$ phân tử trong thể tích $0,02\,m^3$. Tính mật độ phân tử theo đơn vị $10^{25}\,m^{-3}$.
\shortans{3}
\loigiai{
$n_0=N/V=6\times10^{23}/0,02=3\times10^{25}\,m^{-3}$.
}
\end{ex}

% ===== Câu 125 =====
% ID: L12C2B11-124-TLN
% Mức: TH
\begin{ex}
% ID: L12C2B11-124-TLN
% Mức: TH
Trong dạng “Đặc điểm cấu tạo và chuyển động phân tử ở thể khí”, Một lượng khí có thể tích 4 lít được nén còn 2 lít. Mật độ phân tử tăng bao nhiêu lần?
\shortans{2}
\loigiai{
Số phân tử không đổi nên mật độ tỉ lệ nghịch thể tích: $n_2/n_1=V_1/V_2=2$.
}
\end{ex}

% ===== Câu 131 =====
% ID: L12C2B11-125-TLN
% Mức: VD
\begin{ex}
% ID: L12C2B11-125-TLN
% Mức: VD
Trong dạng “Đặc điểm cấu tạo và chuyển động phân tử ở thể khí”, Khí có $N=1,2\times10^{24}$ phân tử và mật độ $3\times10^{25}\,m^{-3}$. Tính thể tích theo lít.
\shortans{40}
\loigiai{
Động năng tịnh tiến trung bình của phân tử khí lý tưởng tỉ lệ thuận với nhiệt độ tuyệt đối theo công thức $\overline{W_đ} = \frac{3}{2}kT$.
Khi nhiệt độ tăng từ $T_1 = 300 \, \mathrm{K}$ lên $T_2 = 600 \, \mathrm{K}$, tỉ số động năng là $\frac{\overline{W_{đ2}}}{\overline{W_{đ1}}} = \frac{T_2}{T_1} = \frac{600}{300} = 2$.
Vậy động năng tịnh tiến trung bình tăng 2 lần.

Đáp số: 2
}
\end{ex}

% ===== Câu 134 =====
% ID: L12C2B11-126-TLN
% Mức: VD
\begin{ex}
% ID: L12C2B11-126-TLN
% Mức: VD
Trong dạng “Đặc điểm cấu tạo và chuyển động phân tử ở thể khí”, Có 2 mol khí. Lấy $N_A=6,02\times10^{23}$. Số phân tử theo đơn vị $10^{23}$ là bao nhiêu?
\shortans{12,04}
\loigiai{
$N=2N_A=12,04\times10^{23}$.
}
\end{ex}

% ===== Câu 137 =====
% ID: L12C2B11-127-TLN
% Mức: VDC
\begin{ex}
% ID: L12C2B11-127-TLN
% Mức: VDC
Trong dạng “Đặc điểm cấu tạo và chuyển động phân tử ở thể khí”, Khí chiếm 5 lít, giãn nở thành 15 lít. Mật độ phân tử giảm bao nhiêu lần?
\shortans{3}
\loigiai{
$n_1/n_2=V_2/V_1=15/5=3$.
}
\end{ex}

% ===== Câu 139 =====
% ID: L12C2B11-128-TN
% Mức: VD
\begin{ex}
% ID: L12C2B11-128-TN
% Mức: VD
Trong dạng “Giải thích áp suất khí bằng va chạm phân tử”, Ba thông số trạng thái cơ bản của một lượng khí là
\choice
{\True áp suất, thể tích, nhiệt độ}
{khối lượng, trọng lượng, nhiệt lượng}
{vận tốc, gia tốc, thời gian}
{lực, công, công suất}
\loigiai{
Số mol khí $n = 2 \, \mathrm{mol}$, số Avogadro $N_A = 6,02 \times 10^{23} \, \mathrm{phân tử/mol}$.
Số phân tử khí có trong bình là $N = n \cdot N_A = 2 \cdot 6,02 \times 10^{23} = 12,04 \times 10^{23} \, \mathrm{phân tử}$.
Vậy số phân tử theo đơn vị $10^{23}$ là 12,04.

Đáp số: 12,04
}
\end{ex}

% ===== Câu 141 =====
% ID: L12C2B11-129-TN
% Mức: VDC
\begin{ex}
% ID: L12C2B11-129-TN
% Mức: VDC
Trong dạng “Giải thích áp suất khí bằng va chạm phân tử”, Khi nén khí ở nhiệt độ không đổi, mật độ phân tử
\choice
{\True tăng}
{giảm}
{không đổi}
{bằng không}
\loigiai{
Đáp án A. Cùng số phân tử trong thể tích nhỏ hơn nên mật độ phân tử tăng.
}
\end{ex}

% ===== Câu 149 =====
% ID: L12C2B11-130-TN
% Mức: VD
\begin{ex}
% ID: L12C2B11-130-TN
% Mức: VD
Trong dạng “Vận dụng mô hình động học phân tử vào thực tế”, Ba thông số trạng thái cơ bản của một lượng khí là
\choice
{\True áp suất, thể tích, nhiệt độ}
{khối lượng, trọng lượng, nhiệt lượng}
{vận tốc, gia tốc, thời gian}
{lực, công, công suất}
\loigiai{
Đáp án A. Trạng thái lượng khí xác định bởi p, $V$, T.
}
\end{ex}

% ===== Câu 151 =====
% ID: L12C2B11-131-TN
% Mức: VDC
\begin{ex}
% ID: L12C2B11-131-TN
% Mức: VDC
Trong dạng “Vận dụng mô hình động học phân tử vào thực tế”, Khi nén khí ở nhiệt độ không đổi, mật độ phân tử
\choice
{\True tăng}
{giảm}
{không đổi}
{bằng không}
\loigiai{
Đáp án A. Cùng số phân tử trong thể tích nhỏ hơn nên mật độ phân tử tăng.
}
\end{ex}

% ===== Câu 159 =====
% ID: L12C2B11-132-TN
% Mức: VD
\begin{ex}
% ID: L12C2B11-132-TN
% Mức: VD
Trong dạng “Đặc điểm cấu tạo và chuyển động phân tử ở thể khí”, Ba thông số trạng thái cơ bản của một lượng khí là
\choice
{\True áp suất, thể tích, nhiệt độ}
{khối lượng, trọng lượng, nhiệt lượng}
{vận tốc, gia tốc, thời gian}
{lực, công, công suất}
\loigiai{
Đáp án A. Trạng thái lượng khí xác định bởi p, $V$, T.
}
\end{ex}

% ===== Câu 161 =====
% ID: L12C2B11-133-TN
% Mức: VDC
\begin{ex}
% ID: L12C2B11-133-TN
% Mức: VDC
Trong dạng “Đặc điểm cấu tạo và chuyển động phân tử ở thể khí”, Khi nén khí ở nhiệt độ không đổi, mật độ phân tử
\choice
{\True tăng}
{giảm}
{không đổi}
{bằng không}
\loigiai{
Đáp án A. Cùng số phân tử trong thể tích nhỏ hơn nên mật độ phân tử tăng.
}
\end{ex}
